% Options for packages loaded elsewhere
\PassOptionsToPackage{unicode}{hyperref}
\PassOptionsToPackage{hyphens}{url}
\documentclass[
  ignorenonframetext,
]{beamer}
\newif\ifbibliography
\usepackage{pgfpages}
\setbeamertemplate{caption}[numbered]
\setbeamertemplate{caption label separator}{: }
\setbeamercolor{caption name}{fg=normal text.fg}
\beamertemplatenavigationsymbolsempty
% remove section numbering
\setbeamertemplate{part page}{
  \centering
  \begin{beamercolorbox}[sep=16pt,center]{part title}
    \usebeamerfont{part title}\insertpart\par
  \end{beamercolorbox}
}
\setbeamertemplate{section page}{
  \centering
  \begin{beamercolorbox}[sep=12pt,center]{section title}
    \usebeamerfont{section title}\insertsection\par
  \end{beamercolorbox}
}
\setbeamertemplate{subsection page}{
  \centering
  \begin{beamercolorbox}[sep=8pt,center]{subsection title}
    \usebeamerfont{subsection title}\insertsubsection\par
  \end{beamercolorbox}
}
% Prevent slide breaks in the middle of a paragraph
\widowpenalties 1 10000
\raggedbottom
\AtBeginPart{
  \frame{\partpage}
}
\AtBeginSection{
  \ifbibliography
  \else
    \frame{\sectionpage}
  \fi
}
\AtBeginSubsection{
  \frame{\subsectionpage}
}
\usepackage{iftex}
\ifPDFTeX
  \usepackage[T1]{fontenc}
  \usepackage[utf8]{inputenc}
  \usepackage{textcomp} % provide euro and other symbols
\else % if luatex or xetex
  \usepackage{unicode-math} % this also loads fontspec
  \defaultfontfeatures{Scale=MatchLowercase}
  \defaultfontfeatures[\rmfamily]{Ligatures=TeX,Scale=1}
\fi
\usepackage{lmodern}
\ifPDFTeX\else
  % xetex/luatex font selection
\fi
% Use upquote if available, for straight quotes in verbatim environments
\IfFileExists{upquote.sty}{\usepackage{upquote}}{}
\IfFileExists{microtype.sty}{% use microtype if available
  \usepackage[]{microtype}
  \UseMicrotypeSet[protrusion]{basicmath} % disable protrusion for tt fonts
}{}
\makeatletter
\@ifundefined{KOMAClassName}{% if non-KOMA class
  \IfFileExists{parskip.sty}{%
    \usepackage{parskip}
  }{% else
    \setlength{\parindent}{0pt}
    \setlength{\parskip}{6pt plus 2pt minus 1pt}}
}{% if KOMA class
  \KOMAoptions{parskip=half}}
\makeatother
\usepackage{longtable,booktabs,array}
\newcounter{none} % for unnumbered tables
\usepackage{calc} % for calculating minipage widths
\usepackage{caption}
% Make caption package work with longtable
\makeatletter
\def\fnum@table{\tablename~\thetable}
\makeatother
\setlength{\emergencystretch}{3em} % prevent overfull lines
\providecommand{\tightlist}{%
  \setlength{\itemsep}{0pt}\setlength{\parskip}{0pt}}
\usepackage{bookmark}
\IfFileExists{xurl.sty}{\usepackage{xurl}}{} % add URL line breaks if available
\urlstyle{same}
\hypersetup{
  hidelinks,
  pdfcreator={LaTeX via pandoc}}

\author{\texorpdfstring{}{}}
\date{}

\begin{document}

\section{Test matrix, mutation testing, and benchmark
suite}\label{test-matrix-mutation-testing-and-benchmark-suite}

\begin{frame}[fragile]{Test matrix and coverage}
\protect\phantomsection\label{test-matrix-and-coverage}
The v0.5.0 Python implementation ships a test suite that, on the
canonical \texttt{uv\ run\ pytest\ tests/\ -\/-cov=py/cogant} run,
reports \textbf{2129 passing} tests with \textbf{86 skips} for optional
dependencies (Rust toolchain, \texttt{matplotlib}, \texttt{tree-sitter}
language grammars, PNG rasterization), plus \textbf{2 expected
\texttt{xfail}} and \textbf{1 \texttt{xpass}} case. End-to-end runtime
is on the order of four minutes on a 2024-class Apple-silicon
workstation (\textbf{238} s in the canonical run); the overall line
coverage of \texttt{py/cogant/} is \textbf{83.42\%} on that run,
measured across \textbf{56628} lines in \textbf{179} source files (see
\texttt{METRICS.yaml}).

\textbf{Table 8. Python interpreter matrix.}

{\def\LTcaptype{none} % do not increment counter
\begin{longtable}[]{@{}
  >{\raggedright\arraybackslash}p{(\linewidth - 4\tabcolsep) * \real{0.3333}}
  >{\raggedright\arraybackslash}p{(\linewidth - 4\tabcolsep) * \real{0.3333}}
  >{\raggedright\arraybackslash}p{(\linewidth - 4\tabcolsep) * \real{0.3333}}@{}}
\toprule\noalign{}
\begin{minipage}[b]{\linewidth}\raggedright
Python version
\end{minipage} & \begin{minipage}[b]{\linewidth}\raggedright
\texttt{pyproject.toml} classifier
\end{minipage} & \begin{minipage}[b]{\linewidth}\raggedright
Status
\end{minipage} \\
\midrule\noalign{}
\endhead
3.11 & \texttt{Programming\ Language\ ::\ Python\ ::\ 3.11} & supported
(minimum version, \texttt{requires-python\ =\ "\textgreater{}=3.11"}) \\
3.12 & \texttt{Programming\ Language\ ::\ Python\ ::\ 3.12} & supported
(canonical CI interpreter; benchmark runs use 3.12.11) \\
3.13 & \texttt{Programming\ Language\ ::\ Python\ ::\ 3.13} &
supported \\
\bottomrule\noalign{}
\end{longtable}
}

All three interpreters are listed in the \texttt{classifiers} block of
\href{../cogant/pyproject.toml}{\texttt{../cogant/pyproject.toml}}. The
declared minimum is Python 3.11 so that the pattern-matching front end
in \texttt{cogant.static.parser.PythonASTParser} can use
\texttt{match}/\texttt{case} statements without a compatibility shim,
and the benchmark suite recorded in
\texttt{benchmarks/results/suite\_20260409.md} was executed on CPython
3.12.11 under macOS arm64.

Module-level coverage is concentrated in the layers that the \textbf{6}
packaged fixtures exercise end-to-end. Table 9 records the coverage of
the algorithmic core (translation, state-space compilation, Markov
blanket extraction, GNN matrix construction, and the reverse
synthesizer) --- the modules whose correctness is load-bearing for every
claim in the manuscript. Numbers are taken from the \texttt{TOTAL}-line
breakdown of the \texttt{uv\ run\ pytest\ -\/-cov} run that produced the
2129/86 pass/skip summary.

\textbf{Table 9. Line coverage of load-bearing modules (canonical v0.5.0
run, 2026-04-10).}

{\def\LTcaptype{none} % do not increment counter
\begin{longtable}[]{@{}
  >{\raggedright\arraybackslash}p{(\linewidth - 4\tabcolsep) * \real{0.2727}}
  >{\raggedleft\arraybackslash}p{(\linewidth - 4\tabcolsep) * \real{0.3636}}
  >{\raggedleft\arraybackslash}p{(\linewidth - 4\tabcolsep) * \real{0.3636}}@{}}
\toprule\noalign{}
\begin{minipage}[b]{\linewidth}\raggedright
Module
\end{minipage} & \begin{minipage}[b]{\linewidth}\raggedleft
Lines
\end{minipage} & \begin{minipage}[b]{\linewidth}\raggedleft
Coverage
\end{minipage} \\
\midrule\noalign{}
\endhead
\texttt{cogant.translate.engine} & 160 & 90\% \\
\texttt{cogant.translate.rules.structural} & 185 & 93\% \\
\texttt{cogant.translate.rules.semantic} & 216 & 92\% \\
\texttt{cogant.translate.rules.behavioral} & 79 & 98\% \\
\texttt{cogant.translate.rules.control} & 47 & 93\% \\
\texttt{cogant.translate.rules.resilience} & 130 & 95\% \\
\texttt{cogant.translate.confidence} & 98 & 97\% \\
\texttt{cogant.statespace.compiler} & 418 & 91\% \\
\texttt{cogant.statespace.variables} & 182 & 98\% \\
\texttt{cogant.statespace.temporal} & 173 & 81\% \\
\texttt{cogant.markov.blanket} & full coverage (reported in the ``60
files skipped due to complete coverage'' block of the canonical run) &
100\% \\
\texttt{cogant.gnn.matrices} & full coverage (same block) & 100\% \\
\texttt{cogant.static.calls} & 151 & 86\% \\
\texttt{cogant.static.dataflow} & 246 & 84\% \\
\texttt{cogant.static.parser} & 236 & 86\% \\
\texttt{cogant.simulate.free\_energy} & 165 & 99\% \\
\texttt{cogant.simulate.runner} & 252 & 67\% \\
\texttt{cogant.simulate.distributions} & 118 & 95\% \\
\texttt{cogant.scoring.drift} & 215 & 99\% \\
\texttt{cogant.scoring.metrics} & 142 & 99\% \\
\texttt{cogant.validate.integrity} & 136 & 94\% \\
\texttt{cogant.validate.schema\_check} & 115 & 95\% \\
\texttt{cogant.validate.provenance\_check} & 73 & 97\% \\
\bottomrule\noalign{}
\end{longtable}
}

The aggregate project-level coverage reported at the end of the run is
\textbf{83.42\%}; the modules that drag the average down are the
visualisation layer (\texttt{cogant.viz.png\_export},
\texttt{cogant.viz.plots}, \texttt{cogant.viz.mermaid}, with residual
gaps where optional \texttt{matplotlib} and \texttt{plotly} code paths
are skipped) plus a small number of scaffolded plugin or provenance
helpers (for example \texttt{cogant.viz.bundle\_site}, an HTML site
generator that requires the \texttt{jinja2} extra and therefore may not
execute under the default \texttt{uv\ sync} environment). The
algorithmic core --- everything that participates in the round-trip
theorem of §9 --- remains at high coverage on the exercised modules in
Table 9. The \texttt{simulate.distributions} and
\texttt{simulate.free\_energy} modules are reported in that table for
the canonical v0.5.0 run; see \texttt{../cogant/CHANGELOG.md} for
release-cycle deltas.
\end{frame}

\begin{frame}[fragile]{Mutation testing}
\protect\phantomsection\label{mutation-testing}
Mutation testing was performed on the algorithmic core modules
(\texttt{gnn/matrices.py}, \texttt{translate/engine.py},
\texttt{markov/blanket.py}, \texttt{statespace/compiler.py},
\texttt{static/dataflow.py}). The canonical \texttt{mutmut} 3.5.0 runner
was evaluated but rejected: on COGANT's test layout \texttt{mutmut}
reported every one of the 403 auto-generated mutants on
\texttt{matrices.py} as ``no tests'' because its v3 trampoline requires
tests to import the mutated module through the
\texttt{mutants/\textless{}path\textgreater{}} shadow tree, and the
project's \texttt{pytest} configuration does not. Rather than ship a
``no tests'' score, the mutation analysis in
\texttt{../cogant/docs/evaluation/MUTATION\_REPORT.md} is based on a
\textbf{hand-picked set of fifteen semantic mutations} that target the
algorithmic predicates, constants, and loop bounds of the above modules;
each mutation was applied, the relevant \texttt{pytest} subset was
rerun, and the mutation was reverted immediately. This is a more
informative experiment than a green \texttt{mutmut} run because it
documents exactly \emph{which} invariants the tests enforce.

\textbf{Table 10. Hand-curated mutation results on COGANT algorithmic
core.}

{\def\LTcaptype{none} % do not increment counter
\begin{longtable}[]{@{}lrrrr@{}}
\toprule\noalign{}
Module & Mutants tested & Killed & Survived & Mutation score \\
\midrule\noalign{}
\endhead
\texttt{gnn/matrices.py} & 5 & 3 & 2 & 60\% \\
\texttt{translate/engine.py} & 3 & 1 & 2 & 33\% \\
\texttt{markov/blanket.py} & 2 & 1 & 1 & 50\% \\
\texttt{statespace/compiler.py} & 2 & 1 & 1 & 50\% \\
\texttt{static/dataflow.py} & 3 & 3 & 0 & 100\% \\
\textbf{Total} & \textbf{15} & \textbf{10} & \textbf{5} &
\textbf{66.7\%} \\
\bottomrule\noalign{}
\end{longtable}
}

The five surviving mutants are documented individually in
\texttt{../cogant/docs/evaluation/MUTATION\_REPORT.md} §``Surviving
mutants --- action required'' (aversive preference path in
\texttt{compute\_C}, sensory↔active boundary role swap in
\texttt{markov/blanket.py},
\texttt{\textgreater{}=}→\texttt{\textgreater{}} boundary flip in
\texttt{\_map\_confidence}, \texttt{CONFIGURATION} neighbour bias in
\texttt{compute\_D}, and the single-pass fixpoint iteration cap). Three
of the five were closed by hardening tests in the same commit:
\texttt{test\_C\_aversive\_preference\_produces\_negative\_log\_pref}
kills the aversive-preference survivor,
\texttt{test\_boundary\_with\_only\_outgoing\_edge\_is\_active} /
\texttt{test\_boundary\_with\_only\_incoming\_edge\_is\_sensory} kill
the Markov-blanket swap, and
\texttt{test\_map\_confidence\_exact\_boundary\_values} kills the
\texttt{\textgreater{}=}→\texttt{\textgreater{}} family. The remaining
two survivors (CONFIGURATION-bias and single-pass fixpoint) are
documented follow-ups that require non-trivial fixture extensions. The
measured score on the hand-picked set is therefore \textbf{10 killed /
15 total = 66.7\%} before hardening, and the documented target after the
follow-ups is 80\% or better.

The modules with the strongest mutation signal are
\texttt{static/dataflow.py} (3 of 3 killed --- every edge-kind string
mutation is caught by the dataflow-tuple-assertion tests) and the
row/column normalisation paths in \texttt{gnn/matrices.py} (4 of 5
killed --- every arithmetic mutation to \texttt{\_normalize\_row},
\texttt{\_DEFAULT\_DIRECT\_MASS}, the \texttt{compute\_C} sign-flip, the
\texttt{compute\_B} axis swap, and the \texttt{compute\_A} fallback
branch is killed by the \texttt{A\_rows\_sum\_to\_one},
\texttt{B\_columns\_sum\_to\_one\_per\_action}, and
\texttt{A\_concentrates\_mass\_on\_direct\_reads} tests). The modules
that drag the overall score down are those whose tests assert structural
invariants (disjointness, normalisation, shape) but do not pin down the
\emph{direction} or \emph{magnitude} of individual entries.
\end{frame}

\begin{frame}[fragile]{Benchmark suite (shipped)}
\protect\phantomsection\label{benchmark-suite-shipped}
A reproducible benchmark harness lives at
\href{../cogant/benchmarks/bench_suite.py}{\texttt{../cogant/benchmarks/bench\_suite.py}}
and writes its canonical results to
\href{../cogant/benchmarks/results/}{\texttt{../cogant/benchmarks/results/}}.
The most recent run committed to the tree
(\texttt{bench(p1.5):\ Rust\ build\ status\ +\ Python\ vs\ Rust\ benchmark\ results},
then superseded by
\texttt{bench(suite):\ reproducible\ 6-fixture\ benchmark\ harness\ with\ stage\ timing\ +\ memory\ +\ GNN\ stats})
executed each fixture for three iterations on CPython 3.12.11 / macOS
arm64 and recorded per-stage wall-clock time, peak memory, and the final
GNN tensor shapes.

\textbf{Table 11. Benchmark suite results (\texttt{suite\_20260409.md},
three iterations per fixture, CPython 3.12.11).}

{\def\LTcaptype{none} % do not increment counter
\begin{longtable}[]{@{}
  >{\raggedright\arraybackslash}p{(\linewidth - 12\tabcolsep) * \real{0.1111}}
  >{\raggedleft\arraybackslash}p{(\linewidth - 12\tabcolsep) * \real{0.1481}}
  >{\raggedleft\arraybackslash}p{(\linewidth - 12\tabcolsep) * \real{0.1481}}
  >{\raggedleft\arraybackslash}p{(\linewidth - 12\tabcolsep) * \real{0.1481}}
  >{\raggedleft\arraybackslash}p{(\linewidth - 12\tabcolsep) * \real{0.1481}}
  >{\raggedleft\arraybackslash}p{(\linewidth - 12\tabcolsep) * \real{0.1481}}
  >{\raggedleft\arraybackslash}p{(\linewidth - 12\tabcolsep) * \real{0.1481}}@{}}
\toprule\noalign{}
\begin{minipage}[b]{\linewidth}\raggedright
Fixture
\end{minipage} & \begin{minipage}[b]{\linewidth}\raggedleft
Wall-clock median (ms)
\end{minipage} & \begin{minipage}[b]{\linewidth}\raggedleft
Wall-clock p95 (ms)
\end{minipage} & \begin{minipage}[b]{\linewidth}\raggedleft
Nodes
\end{minipage} & \begin{minipage}[b]{\linewidth}\raggedleft
Edges
\end{minipage} & \begin{minipage}[b]{\linewidth}\raggedleft
Mappings
\end{minipage} & \begin{minipage}[b]{\linewidth}\raggedleft
Peak memory (MB)
\end{minipage} \\
\midrule\noalign{}
\endhead
\texttt{calculator} & 32 & 35 & 12 & 25 & 11 & 0.0 \\
\texttt{event\_pipeline} & 36 & 37 & 23 & 36 & 21 & 0.1 \\
\texttt{flask\_mini} & 43 & 45 & 26 & 40 & 25 & 0.3 \\
\texttt{flask\_app} & 86 & 86 & 98 & 154 & 68 & 0.3 \\
\texttt{requests\_lib} & 76 & 77 & 98 & 152 & 55 & 0.7 \\
\texttt{json\_stdlib} & 48 & 49 & 29 & 34 & 19 & 0.0 \\
\bottomrule\noalign{}
\end{longtable}
}

The benchmark harness times the bare translation pipeline
(\texttt{ingest}, \texttt{static}, \texttt{normalize}, \texttt{graph},
\texttt{translate}, \texttt{statespace}), so its wall-clock numbers are
approximately an order of magnitude smaller than the end-to-end
roundtrip times of Table 4: Tables 4 and 5 include validation, GNN
package assembly, Mermaid and PNG rasterization, GraphML and Parquet
serialization, and the HTML dashboard write, none of which are part of
the benchmark hot path. For pure translation, every shipped fixture runs
in under 100 ms and consumes less than a megabyte of peak memory; the
stage breakdown in \texttt{suite\_20260409.md} shows that for the
smaller fixtures the dominant cost is \texttt{ingest} (repository walk +
file hashing), and for the larger fixtures (\texttt{flask\_app},
\texttt{requests\_lib}) the dominant cost shifts to \texttt{graph}
construction where \texttt{CallGraphBuilder} walks every
\texttt{ast.Call} node to produce the CALLS edges recorded in Table 5.

Approximate stage breakdown from the same run: \texttt{ingest} 25--30 ms
across all fixtures; \texttt{static} 1--4 ms; \texttt{normalize} 0--3
ms; \texttt{graph} 4--43 ms (dominated by \texttt{flask\_app});
\texttt{translate} 0--3 ms; \texttt{statespace} 0--1 ms. The benchmark
results file also records the per-fixture GNN tensor shapes --- for
example \texttt{flask\_app} produces
\(A \in \mathbb{R}^{21 \times 10}\),
\(B \in \mathbb{R}^{10 \times 10 \times 31}\),
\(C \in \mathbb{R}^{21}\), \(D \in \mathbb{R}^{10}\) --- which match the
state-space compiler outputs of Table 6 up to the benchmark's
independent re-run sampling variance.
\end{frame}

\end{document}
